\section{Datengrundlage und Anwendungsrahmen}
\label{sec:datengrundlage}

\subsection{Alurohr als kontrollierter Testdatensatz}

Die Projektarbeit verwendet eine etwa 48 Sekunden lange Videosequenz eines
Alurohrs als kontrollierten Labordatensatz. Das Rohr repräsentiert ein dünnes,
lineares Objekt, dessen Bildfläche sich mit Kameraposition, Blickwinkel und
Auflösung deutlich verändert. Der Datensatz dient nicht als maßstäblicher
Ersatz für ein im Baugraben liegendes Kabel. Er ermöglicht jedoch, die
Datenübergaben und Fehlerfälle einer vollständigen Rekonstruktionskette unter
wiederholbaren Bedingungen zu untersuchen.

\subsection{Aufnahmekamera und Videoeinstellungen}

Die Videosequenz wurde durchgehend mit der Hauptkamera (1x-Weitwinkelkamera)
des Google Pixel 8 Pro aufgenommen. Diese Kamera besitzt laut
Herstellerangaben einen 50-MP-Sensor mit einer Blende von $f/1{,}68$, einem
Sichtfeld von etwa $82^\circ$ und einer Sensorgröße von $1/1{,}31$ Zoll;
Google weist für die Weitwinkelkamera außerdem optische und elektronische
Bildstabilisierung aus \parencite{googlePixel8ProSpecs}. Die Aufnahme erfolgte
mit der OpenCamera-App. OpenCamera unterstützt eine Belichtungssperre sowie
die Auswahl der Videoauflösung; die automatische Bildausrichtung (Auto-Level)
ist eine separate, optionale Funktion \parencite{openCameraHelp}.

Für den verwendeten Datensatz waren die Belichtung fest eingestellt, die
Auflösung auf $1920\times1080$ Pixel (Full HD) gesetzt und die automatische
Bildlagenkorrektur beziehungsweise nachträgliche Bildausrichtung deaktiviert.
Die Rohdatei \path{data/01_raw/Alurohr_THWS.mp4} ist ein H.264-Video mit
$1920\times1080$ Pixeln, nominal 30 FPS, 1441 Bildern und einer Dauer von rund
48 Sekunden. Die späteren 5- beziehungsweise 2-FPS-
Profile sowie die Zielauflösungen 720p, qHD und Low sind daraus abgeleitete
Pipeline-Arbeitsprofile und keine zusätzlichen Kameraaufnahmen.

Die Bezeichnung ``unverfälschtes Video'' ist dabei auf die fehlende zusätzliche
App-seitige Bildausrichtung und die festgehaltene Belichtung zu beziehen. Die
Aufnahme ist kein Rohsensordatensatz: Das Pixel verarbeitet das Signal bereits
über seine interne Bildverarbeitung und speichert hier ein H.264-Video.
Insbesondere lässt sich aus der Datei allein nicht sicher nachweisen, ob die
vom Gerät bereitgestellte optische beziehungsweise elektronische
Videostabilisierung vollständig deaktiviert war. Die PA behauptet daher keine
vollständig unverarbeitete Sensoraufnahme, sondern dokumentiert ein unter
festgelegten Aufnahmeparametern erzeugtes Rohvideo für die nachgelagerte
Pipeline.

Der Segmentierungsprompt lautet \texttt{pipe}. Die in der Projektarbeit
untersuchten Profilvarianten aus zwei zeitlichen Abtastungen (2 und 5 FPS) und
drei räumlichen Auflösungen ($1280\times720$, $960\times540$ und
$640\times360$ Pixel) sind in Kapitel~\ref{sec:versuchsaufbau} zusammengefasst;
für den Datensatz genügt hier der Hinweis, dass ein Lauf bei 5 FPS rund 240
Frames und bei 2 FPS rund 96 Frames umfasst. Ein Lauf verwendet in allen
Verarbeitungsstufen denselben Frame-Satz, sodass Bildnamen, Masken, Kameras
und feste Evaluationsansichten innerhalb eines Laufs konsistent bleiben.

\subsection{Scope des linearen Einobjektfalls}

Die Centerline-Nachverarbeitung ist bewusst auf ein einzelnes lineares Objekt
begrenzt. Verzweigte Leitungsnetze erfordern eine zusätzliche Trennung echter
Abzweigungen von Skelettartefakten. Gerade auf dünnwandigen und verrauschten
Meshes kann die Skelettierung Medialflächen\footnote{Medialflächen sind
zweidimensionale, flächige Skelettbereiche: Sie entstehen, wenn von einem
inneren Punkt mehrere Punkte der Objektoberfläche gleich weit entfernt sind
(zum Beispiel bei flachen oder dünnwandigen Strukturen). Für eine
Linientrasse sind sie nicht direkt verwertbar und müssen auf ihren
Durchmesserpfad reduziert werden.\label{fn:medialflaechen}} und zahlreiche
kurze Stachel-Äste erzeugen. Der produktive \texttt{single}-Modus reduziert
deshalb die größte zusammenhängende Skelettkomponente auf ihren
Durchmesserpfad.

Das primäre analytische Produkt der PA ist die lokale Centerline. Masken,
COLMAP-Modell, STS-Gaussians und Mesh sind notwendige und diagnostisch wichtige
Zwischenprodukte. CSV-, B-Spline- und GeoJSON-Dateien werden ebenfalls erzeugt,
aber wegen der noch fehlenden unabhängigen Transformations- und
Vermessungsreferenz nicht als global genaue Endprodukte bewertet.

\subsection{Referenzlage und Aussagegrenzen}

Für den aktuellen Testdatensatz liegt keine unabhängig eingemessene
GNSS-Referenzkurve vor. Vorhandene GCP-Dateien bilden nur einen synthetischen
beziehungsweise lokalen Test der Dateiverträge. Fehlt eine validierte
4x4-Ähnlichkeitstransformation, kann die Pipeline einen explizit bezeichneten
Translation-Fallback ausführen. Dieser prüft den Exportpfad, aber weder Rotation
und Maßstab noch Lage- und Höhengenauigkeit.

Daraus folgt die zentrale Abgrenzung: Die PA bewertet Funktionsfähigkeit,
Robustheit, Reproduzierbarkeit, objektmaskierte Ansichtsqualität und relative
Unterschiede zwischen Meshrouten. Ein Nachweis der später angestrebten
Toleranz von $\pm10\,\mathrm{cm}$ ist erst mit realen GCP-/GNSS-Messungen,
einer validierten Transformation und geometrischen Vergleichsmetriken möglich.
